\documentclass[../numerical_methods.tex]{subfiles}
\begin{document}
\section{Aula 06 - 19 de Março, 2026}
\subsection{Motivações}
\begin{itemize}
	\item Exercícios de Numpy.
\end{itemize}
\subsection{Exercícios de Numpy}
Recomenda-se fortemente que os exercícios sejam feitos \textit{sem consultar as respostas!} Além disso, para todos os códigos, presume-se que foi colocado o comando \texttt{import numpy as np}.
\subsubsection{Exercício 1}
Gere a lista de todos os números inteiros entre 0 e 19. Converta a lista em um numpy array bidimensional (matriz) com 4 linhas e 5 colunas.

\begin{minted}[bgcolor=DarkBackground, linenos=true]{Python}
#| colab: {base_uri: https://localhost:8080/}
l = [i for i in range(20)]
A = np.array(l).reshape(4,5)
print(A)
\end{minted}


\subsubsection{Exercício 2}
Construa um array unidimensional com 30 números inteiros escolhidos de forma randômica no intervalo entre 0 e 20. Reformate o array para que se torne uma matriz `A` com 5 linhas e 6 colunas.

\begin{minted}[bgcolor=DarkBackground, linenos=true]{Python}
A = np.random.randint(0,20,30).reshape(5,6)
print(A)
\end{minted}


\subsubsection{Exercício 3}
Crie uma view da matriz `A` gerada na célula anterior contendo:
- as linhas de `A` com índices 0,1 e 3
- as linhas de `A` com índice 1 e 2 e as colunas com índice 0,2 e 4
- as linhas de `A` com índice 1 e 3 e as colunas com índice 1 e 3

\begin{minted}[bgcolor=DarkBackground, linenos=true, breaklines]{Python}
print('Matriz A:')
print(A,'\n')

A0 = A[[0,1,3]] # linhas de `A` com índices 0,1 e 3
print("linhas de `A` com indices 0,1 e 3")
print(A0,'\n')

A1 = A[1:3,[0,2,4]] # linhas de `A` com índice 1 e 2 e as colunas
#com índice 0,2 e 4
print("linhas de `A` com indice 1 e 2 e colunas com indice 0,2 e 4")
print(A1,'\n')

A2 = A[[1,3]][:,[1,3]] #linhas de `A` com índice 
#1 e 3 e as colunas com índice 1 e 3
print("linhas de `A` com indice 1 e 3 e colunas com indice 1 e 3")
print(A2)
\end{minted}


\subsubsection{Exercício 4}
Construa um array `A` bidimensional 5x6 com 30 números inteiros escolhidos de forma randômica no intervalo entre 0 e 20.
- faça uma cópia do array `A` criado, chamando a cópia de `Ac`
- Substitua todos os elementos do array `A` que sejam maiores que 10 pelo valor -1
- Substitua todos os elementos do array `Ac` que sejam maiores que 7 e menores que 15 pelo valor -1

Dica: Gere máscaras booleanas para realizar as modificações nos arrays.

\begin{minted}[bgcolor=DarkBackground, linenos=true, breaklines]{Python}
A = np.random.randint(0,20,30).reshape(5,6)
print('A original')
print(A)

Ac = np.copy(A)  # copia de A

maskA = (A > 10)
A[maskA]  = -1
print('\nA modificada com -1 para elementos maiores que 10')
print(A)

maskAc = ((Ac > 7) & (Ac < 15))
Ac[maskAc] = -1
print('\nAc modificada com -1 para elementos maiores que 7 e menores que 15')
print(Ac)
\end{minted}


\subsubsection{Exercício 5}
Construa um array `A` com 20 linhas e 10 colunas onde os elementos são os números interios de  1 até 200.
- Crie uma view de `A` chamada `A\_lpares` contendo apenas as linhas de `A` com índice par.
- Crie uma view de `A` chamada `A\_lpares\_cimpares` contendo as linhas `A` com índice par e colunas com índice ímpar.

\begin{minted}[bgcolor=DarkBackground, linenos=true]{Python}
A = np.arange(1,201).reshape(20,10)

A_lpares = A[np.arange(0,A.shape[0],2)]
print('A_lpares')
print(A_lpares)

print('\nA_lpares_cimpares')
A_lpares_cimpares = A[np.arange(0,A.shape[0],2)]
[:,np.arange(1,A.shape[1],2)]
print(A_lpares_cimpares)
\end{minted}


\subsubsection{Exercício 6}
Substitua os elementos do array `A` do exercício anterior por -1 quando pelo menos um dos índices do elemento é par.

Dica: Construa uma máscara booleana

\begin{minted}[bgcolor=DarkBackground, linenos=true, breaklines]{Python}
mask = [(i%2==0) or (j%2==0) for i in range(A.shape[0]) for j in range(A.shape[1])]
mask = np.array(mask).reshape(A.shape[0],A.shape[1])
A[mask] = -1

print(A)
\end{minted}


\subsubsection{Exercício 7}
Considere o array `X` de números tipo float gerado a partir de uma distribuição uniforme no intervalo entre 0 e 1. Escreva um código para encontrar os três menores valores do array.

Dica: Converta o array em uma lista e ordene.

\begin{minted}[bgcolor=DarkBackground, linenos=true]{Python}
X = np.random.uniform(0,1,30).reshape(6,5)
print(X)

X_sorted = sorted(X.ravel())
print('Os tres menores valores: ',X_sorted[0:3])
\end{minted}


\subsubsection{Exercício 8}
Escreva uma função chamada `troca\_colunas` que recebe um array e o índice de duas colunas como parâmetros e retorna o array com as colunas trocadas. Por exemplo:
\begin{minted}[bgcolor=DarkBackground, linenos=true]{Python}
X = np.array([[0, 1, 2, 3, 4, 5],
		[0, 1, 2, 3, 4, 5],
		[0, 1, 2, 3, 4, 5]])

X = troca_colunas(X,1,3)
print(X)
\end{minted}

deve resultar em
\begin{minted}[bgcolor=DarkBackground, linenos=true]{Python}
	[[0 3 2 1 4 5]
			[0 3 2 1 4 5]
			[0 3 2 1 4 5]]
\end{minted}


\begin{minted}[bgcolor=DarkBackground, linenos=true, breaklines]{Python}
def troca_colunas(A,i,j):
    aux = np.copy(A[:,i])
    A[:,i] = A[:,j]
    A[:,j] = aux
    return(A)

X = np.array([[0, 1, 2, 3, 4, 5], [0, 1, 2, 3, 4, 5], [0, 1, 2, 3, 4, 5]])
print(X)

X = troca_colunas(X,1,3)
print(X)
\end{minted}


\subsubsection{Exercício 9}
Os métodos vstack e hstack do numpy empilha matrizes na vertical e horizontal, respectivamente. Por exemplo:
\begin{minted}[bgcolor=DarkBackground, linenos=true]{Python}
A = [[1.0, 1.0, 1.0], [1.0, 1.0, 1.0]]
B = [[2.0, 2.0, 2.0], [2.0, 2.0, 2.0]]

C = np.vstack((A,B))
print(C)
\end{minted}

resulta em
\begin{minted}[bgcolor=DarkBackground, linenos=true]{Python}
	[[1. 1. 1.]
			[1. 1. 1.]
			[2. 2. 2.]
			[2. 2. 2.]]
\end{minted}

Escreva um código para gerar duas matrizes `A` e `B` de números inteiros randômicos no intervalo entre 0 e 10, onde:
- A matriz `A` é 5x6
- A matriz `B` é 5x3

Utilize o método hstack para empilhar as matrizes

O que acontece se tentar utilizar o método vstack? Explique o motivo da mensagem gerada.

\begin{minted}[bgcolor=DarkBackground, linenos=true, breaklines]{Python}
A = np.random.randint(0,10,30).reshape(5,6)
B = -1*np.random.randint(0,10,15).reshape(5,3)
print(A)
print(B)

C = np.hstack((A,B))
print(C)

# o método vstack nao pode ser utilizado pois o número de colunas
# não é o mesmo nas duas matrizes, somente o número de linhas coincide
\end{minted}


\subsubsection{Exercício 10}
Escreva um código para, data uma matriz, zerar todos os elementos abaixo da diagonal principal da matriz. Por exemplo:
\begin{minted}[bgcolor=DarkBackground, linenos=true]{Python}
	[[2 8 3 4 8]
			[4 8 7 7 3]
			[8 2 3 1 4]
			[8 6 3 8 1]
			[6 4 4 6 1]]
\end{minted}

deve resultar em
\begin{minted}[bgcolor=DarkBackground, linenos=true]{Python}
	[[2 8 3 4 8]
			[0 8 7 7 3]
			[0 0 3 1 4]
			[0 0 0 8 1]
			[0 0 0 0 1]]
\end{minted}

\begin{minted}[bgcolor=DarkBackground, linenos=true, breaklines]{Python}
A = np.random.randint(1,10,25).reshape(5,5)

mask = [True if i>j else False for i in range(A.shape[0]) for j in range(A.shape[1])]
mask = np.array(mask).reshape(A.shape[0],A.shape[1])
print(mask)

A[mask] = 0
print(A)

# uma alternativa é utilizar o método 'triu' do numpy
A = np.random.randint(1,10,25).reshape(5,5)
print(A)
A = np.triu(A)
print(A)
\end{minted}

\end{document}
